\section{Finite relational models}\label{section: appendix: rel}

In this appendix we prove the relational part of \Cref{theorem: main}.
\begin{theorem}\label{theorem: main REL}
    Let $h \in \nat$ and let $\term$ be a \kl{term} (over any variables) with $\sh(\term) \le h$.
    Then there is a finite relational model in which $\Qt{h}$ and $\term$ differ. In particular, $\REL \not\models \Qt{h} = \term$.
\end{theorem}
The structure of the proof is the same as for $\RSUB$; we use finite cyclic versions of the \kl{chain models} and describe only the differences.
The cyclic setting is in fact simpler in the Kleene star case, since closed-walk weights then form a subgroup of a finite cyclic group.

\subsection{Cyclic models}
\AP Let $N$ be even and $a \in \pnat$.
The \intro{cyclic model} $\intro*\ModR{a}$ (with respect to $N$) is the \kl{relational valuation} on the finite set $\intro*\chainmC{N} \defeq \inte_N \times \range{1, m}$ with
\begin{align*}
    z &\mapsto \set{(\pt{t}{i}, \pt{t+k}{i}) \mid k \in \range{0, N/2 - 1}}, &
    x &\mapsto \set{(\pt{t}{i}, \pt{t + a + p i}{i})}, \\
    c &\mapsto \set{(\pt{t}{i}, \pt{t+1}{j})}, &
    y &\mapsto \set{(\pt{t}{i}, \pt{t + p i - 2}{i})},
\end{align*}
where times are taken modulo $N$, and all other variables are mapped to $\emptyset$; we use the relational semantics $\semR{\bl}{}$.
\AP We call $W \defeq \range{0, N/2 - 1}$ the \intro{window}: $z$ relates two points of the same \kl{phase} whose time difference lies in $W$.
The rotation $\pt{t}{i} \mapsto \pt{t + s}{i}$ is an automorphism, so with the \kl{difference sets}
$\intro*\dsetR{\term}{i}{j}(a) \defeq \set{n \in \inte_N \mid (\pt{0}{i}, \pt{n}{j}) \in \semR{\term}{\ModR{a}}}$
we have $(\pt{t}{i}, \pt{t'}{j}) \in \semR{\term}{\ModR{a}} \iff t' - t \in \dsetR{\term}{i}{j}(a)$.
The composition formula of \Cref{section: upper bound proof} and~\eqref{equation: star walks} hold verbatim, with sums taken in $\inte_N$, and complements are taken in $\inte_N$.

\subsection{The witnesses in cyclic models}
\AP We say that $N$ is \intro{admissible} for $(a, h)$ if $\fseq{r}(a) + p i$ divides $N$ and $2 \fseq{r}(a) + 2 + p m < N/2$ for all $r \le h$ and $i \in \range{1, m}$.
\begin{lemma}\label{lemma: one step cyclic}
    Let $L \defeq \lcm(q_1, \dots, q_m)$, and assume that $L + p i$ divides $N$ for all $i$ and that $2L + 2 + p m < N/2$.
    If $\dsetR{Q}{i}{i}(a) \cap W = q_i \nat \cap W$ for all $i$, then $\dsetR{\Phit(Q)}{i}{i}(a) \cap W = (L + p i) \nat \cap W$ for all $i$.
    Consequently, if $N$ is \kl{admissible} for $(a, h)$, then $\dsetR{\Qt{h}}{1}{1}(a) \cap W = \qseq{h}(a) \nat \cap W$.
\end{lemma}
\begin{proof}
    We follow the proof of \Cref{lemma: one step}.
    The \kl{term} $\Zthree$ relates the pairs in the same \kl{phase} with time difference in $\range{3, N/2 - 1}$, and $z \inter \com{Q}$ only depends on $Q$ inside the \kl{window}.
    For $n \in \range{3, N/2 - 1}$, the pair $(\pt{t}{i}, \pt{t + n}{i})$ is in $c \compo (z \inter \com{Q}) \compo c$ iff $L \ndiv n - 2$, since $n - 2 \in W$.
    So $\Mt(Q)$ relates exactly the pairs $(\pt{t}{i}, \pt{t + 2 + r L}{i})$ with $r \ge 1$ and $2 + rL < N/2$.
    The \kl{term} $c \compo z$ relates exactly the pairs with time difference in $\range{1, N/2}$.
    For two $\Mt(Q)$-successors $v = \pt{t + 2 + rL}{i}$ and $v' = \pt{t + 2 + r'L}{i}$ of $\pt{t}{i}$, the time difference from $v'$ to $v$ is $(r - r')L \in \range{1, N/2}$ if $r' < r$, and $N - (r' - r)L > N/2$ if $r' > r$, as $(r' - r)L < N/2$.
    Hence $\Jt(Q)$ relates exactly the pairs $(\pt{t}{i}, \pt{t + L + 2}{i})$, and $\Jt(Q) \compo y$ those with time difference $L + p i$ in \kl{phase} $i$.
    Its reflexive transitive closure relates time differences in the subgroup of $\inte_N$ generated by $L + p i$, which consists of the multiples of $L + pi$ as $L + p i$ divides $N$.
    The last claim follows by induction as in \Cref{corollary: witness periods}, starting from $\Qt{0} = x^{*}$ (note that $a + p i$ divides $N$).
\end{proof}

\subsection{Periodicity relative to the boundary}
\AP Let $\intro*\bnd \defeq \set{0, N/2} \subseteq \inte_N$, a subgroup, and let $\intro*\far(n)$ be the distance from $n \in \inte_N$ to $\bnd$ along the cycle; $\far$ is invariant under translation by $\bnd$.
\AP For $n \in \inte_N$ let $\intro*\sgn(n) \in \inte$ be the signed distance from $n$ to the nearest point of $\bnd$, so that $n \in \sgn(n) + \bnd$.
We say that $X \subseteq \inte_N$ is \kl{$P$-periodic from $T$} if $n \in X \Leftrightarrow n + P \in X$ for every $n$ such that $\far(n + s) \ge T$ for all $s \in \range{0, P}$; that is, $X$ is $P$-periodic on every stretch that keeps distance at least $T$ from $\bnd$.
Complements, multiples of periods, and larger \kl{preperiods} behave as before, and every subset of $\set{n \mid \far(n) < T}$ is $P$-periodic from $T$.
\begin{fact}\label{fact: periods cyclic}
    Let $X, Y \subseteq \inte_N$.
    \begin{enumerate}
        \item\label{fact: periods cyclic: min} Let $X$ be $P$-periodic from $T$, and let $J = \range{l, l + q + P}$ with $\far(n) \ge T$ for all $n \in J$. If $n \in X \Leftrightarrow q \mid n$ for all $n \in J$, then $q \mid P$.
        \item\label{fact: periods cyclic: sum} If $X$ and $Y$ are $P$-periodic from $T_X$ and $T_Y$, respectively, then $X + Y$ is $P$-periodic from $T_X + T_Y + 2P$.
    \end{enumerate}
\end{fact}
\begin{proof}
    (\ref{fact: periods cyclic: min}) As in \Cref{fact: periods}.
    (\ref{fact: periods cyclic: sum})
    Let $T \defeq T_X + T_Y + 2P$, and let the stretch $\range{n, n + P}$ keep distance at least $T$ from $\bnd$.
    Every point $n + u - \varepsilon$ with $u \in \range{-P, 2P}$ and $\abs{\varepsilon} < T_X$ is within distance $P + T_X$ of this stretch, so it has $\far \ge T_Y + P$.
    ($\Rightarrow$) Let $n = \alpha + \beta$ with $\alpha \in X$, $\beta \in Y$.
    If the stretch $\range{\alpha, \alpha + P}$ keeps distance at least $T_X$ from $\bnd$, then $\alpha + P \in X$ and $n + P = (\alpha + P) + \beta$.
    Otherwise $\alpha + s = \kappa + \varepsilon$ for some $s \in \range{0, P}$, $\kappa \in \bnd$, and $\abs{\varepsilon} < T_X$.
    Then $\beta + u = n + s + u - \varepsilon - \kappa$, and since $\far$ is invariant under $\bnd$ (here we use that $\bnd$ is a subgroup), every point of $\range{\beta, \beta + P}$ has $\far \ge T_Y + P$.
    So $\beta + P \in Y$ and $n + P = \alpha + (\beta + P)$.
    ($\Leftarrow$) Symmetric, using the stretches $\range{\alpha - P, \alpha}$ and $\range{\beta - P, \beta}$.
\end{proof}

\subsection{Upper bound in cyclic models}
\AP We now fix a polynomial $N \in \inte[a]$ taking even positive values, and consider the family $(\ModR{a})_a$ with respect to $N(a)$; we represent elements of $\inte_{N(a)}$ by integers in $\range{0, N(a) - 1}$.
\Cref{claim: tame} makes sense verbatim with this notion of periodicity.
\begin{proposition}\label{proposition: upper bound cyclic}
    If $\deg N > \dg{\term}$, then \Cref{claim: tame} holds for $\term$ in the \kl{cyclic models} with respect to $N$.
\end{proposition}
\begin{proof}
    We follow the induction of \Cref{section: upper bound proof}; the degree condition ensures that all \kl{preperiods} of subterms are eventually much smaller than $N(a)$.
    Variables and constants have the \kl{difference sets} $\emptyset$, $\set{0}$, $\set{1}$, $\set{p i - 2}$, $\set{a + p i}$, or $W$, all $1$-periodic from $O(a)$.
    Complement, union, and composition are handled as before, using \Cref{fact: periods cyclic} instead of \Cref{fact: periods}; definability is clear, as all sets are subsets of $\range{0, N(a) - 1}$ and reduction modulo $N(a)$ is expressed by $\exists k\,(n = n' - k N(a))$.

    For the Kleene star, let $B \defeq \thr{F} + 2 \per{F}$, and call an element $s$ of an edge set \kl{small} if $\far(s) \le B$.
    As in~\eqref{equation: small elements}, if $\far(s) > B$, write $s = \kappa + \sgn(s)$ with $\kappa \in \bnd$; all points between $\kappa + T$ and $s$ keep distance at least $T$ from $\bnd$, so $s$ is an integer combination, with coefficient sum $1$, of two small elements of the same edge set.
    Hence, as in \Cref{lemma: small cycle}, every element of $\cyc{C}$ is an integer combination of \kl{weights} of closed walks of length at most $m$ with small elements.
    If all of them lie in the subgroup $\bnd$, then $\cyc{C} \subseteq \bnd$.
    Otherwise, one of them has \kl{weight} $w = \sum_e (\kappa_e + \sigma_e) \notin \bnd$ with $\abs{\sigma_e} \le B$, so $0 < \abs{\sgn(w)} \le m B$ (as $m B < N/4$).
    We therefore let $\cyw{C}(a)$ be the least non-zero $\abs{\sgn(w)}$ over \kl{weights} $w$ of closed walks inside $C$ of length at most $m$; it is a polynomial of degree at most $d$ on each refined residue class by \Cref{theorem: BGW}, and we set
    \[\per{F^{*}} \defeq \per{F} \cdot \prod_{C \,:\, \cyc{C} \not\subseteq \bnd} 2 \cyw{C}.\]
    The degree bounds of \Cref{table: invariants} are unaffected by the factor $2$.

    It remains to show that the set of \kl{weights} of walks from $i$ to $j$ is $\per{F^{*}}(a)$-periodic from $O(B^2 + \per{F^{*}}(a))$.
    We decompose walks along components as in the proof of \Cref{lemma: walk periods}, and keep its notation $W_C(e, f)$, $G_e$, $w_0$, $\rho$, and~\eqref{equation: W sandwich}.
    If $\cyc{C} \subseteq \bnd$, then $G_e \subseteq \bnd$, and every element of $W_C(e, f)$ has $\far \le m B$, so $W_C(e, f)$ is periodic from $m B + 1$ with any period.
    Otherwise, $G_e$ is a submonoid of the finite group $\inte_N$, hence a subgroup $g \inte_N$ with $g \mid N$, and~\eqref{equation: W sandwich} gives $W_C(e, f) = w_0 + G_e$.
    Let $w$ be a closed walk \kl{weight} at some $v \in C$ with $\abs{\sgn(w)} = \cyw{C}(a)$.
    Going from $e$ to $v$ and back gives $\pi_v, \pi_v + w \in G_e$, so $w \in G_e$ and $2 w = 2 \sgn(w) \in G_e$, because $2 \bnd = \set{0}$.
    Hence $g$ divides the integer $2\cyw{C}(a)$, and the coset $W_C(e, f)$ is $\per{F^{*}}(a)$-periodic everywhere.
    \Cref{fact: periods cyclic}(\ref{fact: periods cyclic: sum}) combines the at most $2m$ summands.
    (No Frobenius-type argument is needed here.)

    Finally, definability follows as in \Cref{lemma: closure}, but more simply: the edge sets are bounded by $N(a)$, so by \Cref{lemma: simplex} they are finite unions of sets $X_\tau(a)$, and $n$ is a walk \kl{weight} iff $0 \le n < N(a)$ and $n = \sum_\tau n_\tau - y N(a)$ for some $(k_\tau)_\tau$ in the Parikh image, $n_\tau \in \kfold{k_\tau}{X_\tau(a)}$, and $y \in \inte$.
\end{proof}

\subsection{Proof of the theorem}
We choose $m$ and $p$ as in \Cref{section: strategy}, and let
\[N(a) \defeq 2 a^{\dg{\term} + 1} \prod_{r = 0}^{h} \prod_{i = 1}^{m} \left(\Aseq{r}(a) + p i \Dm^{\dseq{r}}\right) \in \inte[a].\]
Since $\Aseq{r}(a) + p i \Dm^{\dseq{r}} = \Dm^{\dseq{r}} (\fseq{r}(a) + p i)$ for $a \in \mathcal{A}$, the number $N(a)$ is \kl{admissible} for $(a, h)$ for all large $a \in \mathcal{A}$.
By \Cref{proposition: upper bound cyclic}, \Cref{claim: tame} holds for $\term$ in the \kl{cyclic models} with respect to $N$.
Fix a residue class $r$ modulo the modulus $M$ that meets $\mathcal{A}$ infinitely often, and let $P$ be the period polynomial and $T(a) = O(a^{\dg{\term}})$ the \kl{preperiod} on it.
Suppose that $\Qt{h}$ and $\term$ agree in all finite relational models, in particular in every $\ModR{a}$.
By \Cref{lemma: one step cyclic}, $\dsetR{\term}{1}{1}(a) \cap W = \qseq{h}(a) \nat \cap W$.
The interval $J \defeq \range{T(a), T(a) + \qseq{h}(a) + P(a)}$ lies in $W$ and keeps distance at least $T(a)$ from $\bnd$ for large $a$, because $N(a)/2 \ge a^{\dg{\term} + 1} \qseq{h}(a)$.
So $\qseq{h}(a) \mid P(a)$ by \Cref{fact: periods cyclic}(\ref{fact: periods cyclic: min}) for all large $a \in \mathcal{A}$ with $a \equiv r \pmod M$, and the B\'ezout argument of \Cref{section: strategy} yields a contradiction.
\qed
